\chapter{PROCEDIMENTO EXPERIMENTAL}
\label{cap:procedimento_experimental}

Este capítulo descreve o procedimento experimental utilizado para comparar as estratégias de alocação de Markowitz e Black--Litterman em carteiras de ações negociadas na B3. O experimento utiliza séries históricas de mercado para estimar retornos com o XGBoost e manchetes financeiras para incorporar informações de sentimento às opiniões do modelo Black--Litterman. Em um experimento complementar, são comparados diferentes modelos de previsão de retorno utilizando os mesmos ativos, horizontes e datas de avaliação.

São apresentados a base de dados, as etapas de pré-processamento,
a configuração dos modelos e as regras da simulação financeira,
realizada ao longo dos pregões de 2025. O capítulo também define
as métricas utilizadas para avaliar as previsões, a classificação
de sentimento e o desempenho das carteiras. O código desenvolvido
neste trabalho está disponível publicamente%
\footnote{\url{https://github.com/igxxtl/toledo-portfolio-optimizationn}}.

\section{Base de Dados}
\label{sec:base_dados}

\subsection{Universo de ativos e fontes quantitativas}

O universo do experimento contém vinte ativos negociados
na B3, selecionados para contemplar diferentes setores
e apresentar disponibilidade de dados históricos de
mercado e de manchetes jornalísticas. Os códigos
considerados são ABEV3, B3SA3, BBAS3, BBDC4, CMIG4, EQTL3,
GGBR4, ITUB4, LREN3, MGLU3, PETR3, PETR4, PRIO3, RADL3,
RAIL3, SANB11, SBSP3, SUZB3, VALE3 e WEGE3. A amostra inclui companhias dos setores financeiro, elétrico, industrial, de consumo, infraestrutura e \emph{commodities}. PETR3 e PETR4 possuem séries de preços distintas, mas compartilham informações textuais por pertencerem à mesma empresa. As carteiras avaliadas contêm de uma a vinte ações desse universo.

As séries financeiras foram obtidas da Yahoo Finance
por meio da biblioteca \texttt{yfinance} e armazenadas
em arquivos locais. São utilizados registros diários
de abertura, máxima, mínima, fechamento e volume.
O histórico disponível inicia-se em janeiro de 2018.
Os cálculos utilizam os preços de fechamento registrados
nesses arquivos.

\newpage Na construção dos atributos candidatos do XGBoost,
utilizam-se duas séries exógenas ao histórico de cada ação:
o Ibovespa, que representa o comportamento do mercado
acionário brasileiro, e a taxa de câmbio entre o dólar
americano e o real (USD/BRL), que fornece informação
macroeconômica complementar. A seleção dos atributos
utilizados pelo estimador é descrita na
Seção~\ref{sec:estimador_xgb}.

\subsection{Notícias financeiras e períodos de avaliação}

A base textual reúne manchetes recuperadas pelo serviço
Google News, utilizando sua interface RSS, com idioma
português do Brasil e localização brasileira. As consultas
consideram os códigos de negociação e os nomes das
empresas, abrangendo janeiro de 2023 a dezembro de 2025.
A unidade de análise é a manchete, o corpo das reportagens
não participa da classificação.

A coleta bruta contém 7.835 registros. Após deduplicação
e filtragem de pertinência, a base classificada contém
4.681 manchetes: 2.418 positivas, 1.253 negativas e
1.010 neutras. A existência de uma notícia e a habilitação
de uma opinião são condições distintas, pois a configuração
de alocação exige pelo menos três manchetes no dia
elegível.

O ano de 2024 fornece os dados de validação quantitativa
e a avaliação financeira corresponde aos 250 pregões
entre 2 de janeiro e 30 de dezembro de 2025. A
~\ref{tab:periodos_experimentais} resume os períodos
utilizados.

\begin{table}[H]
\centering
\caption{Períodos de referência do experimento.}
\label{tab:periodos_experimentais}
{\footnotesize
\begin{tabular}{p{0.29\linewidth}p{0.25\linewidth}p{0.35\linewidth}}
\hline
Etapa & Período & Finalidade \\ \hline
Histórico financeiro &
Desde janeiro de 2018 &
Construção de atributos e treinamento quantitativo \\

Notícias financeiras &
Janeiro de 2023 a dezembro de 2025 &
Classificação e agregação de sentimento \\

Validação quantitativa &
Dois semestres de 2024 &
Avaliação dos atributos e definição dos horizontes \\

Previsões de avaliação &
Quatro trimestres de 2025 &
Geração de previsões por atualização temporal \\

Simulação financeira &
250 pregões de 2025 &
Comparação das estratégias de investimento \\

Estimação de risco &
126 pregões até cada sinal &
Cálculo da covariância, do prior e da média histórica \\
\hline
\end{tabular}}
\par
Fonte: Elaborado pelo autor
\end{table}

A ~\ref{fig:divisao_temporal} apresenta a divisão
temporal do experimento. A validação utiliza os dois
semestres de 2024, e o teste abrange os quatro trimestres
de 2025. A cada janela, o modelo é retreinado com os
dados anteriores ao seu início, utilizando apenas
observações cujo retorno-alvo já era conhecido.

\begin{figure}[htbp]
    \centering
    \caption{Divisão temporal do treinamento, da validação
    e do teste do XGBoost.}
    \label{fig:divisao_temporal}

    \begin{tikzpicture}[
        >=Stealth,
        node distance=0.6cm,
        etapa/.style={
            draw,
            thick,
            rounded corners=4pt,
            align=center,
            text width=0.22\linewidth,
            minimum height=1.7cm,
            inner sep=4pt,
            font=\small
        }
    ]
        \node[etapa, fill=gray!15] (treino)
            {\textbf{Treino inicial}\\2018--2023};

        \node[etapa, fill=blue!10, right=of treino] (validacao)
            {\textbf{Validação}\\2024\\Janelas semestrais};

        \node[etapa, fill=green!10, right=of validacao] (teste)
            {\textbf{Teste}\\2025\\Janelas trimestrais};

        \draw[->, thick] (treino) -- (validacao);
        \draw[->, thick] (validacao) -- (teste);
    \end{tikzpicture}

    \par\smallskip
    Fonte: Elaborado pelo autor
\end{figure}

\section{Pré-processamento de Dados}
\label{sec:pre_processamento}

\subsection{Séries financeiras e atributos quantitativos}

Os registros financeiros são ordenados cronologicamente
e verificados quanto à existência de datas duplicadas,
preços de fechamento inválidos e inconsistências.
As séries são alinhadas pelas datas de pregão. Na
simulação financeira, são utilizadas as datas com
fechamento disponível para todas as ações do universo.

Seja $P_{i,t}$ o preço de fechamento da ação $i$ no
pregão $t$. O log-retorno diário é definido por
\begin{equation}
r_{i,t}
=
\ln\left(\frac{P_{i,t}}{P_{i,t-1}}\right).
\label{eq:retornos_cap4}
\end{equation}

Utiliza-se o log-retorno porque ele é aditivo no tempo:
a soma dos log-retornos diários ao longo de um horizonte
equivale ao logaritmo da razão entre os preços de
fechamento final e inicial. Essa propriedade permite
expressar os retornos acumulados nos horizontes de
5, 10 e 21 pregões utilizados como alvos do XGBoost.

Os atributos quantitativos representam o comportamento
histórico da ação e sua relação com o mercado. São
considerados retornos defasados, momentum, volatilidade,
volume, médias móveis, indicadores técnicos e medidas
relativas ao Ibovespa. Os conjuntos candidatos também
incluem atributos cambiais e medidas do comportamento
conjunto das ações.

Entre os atributos relativos ao mercado, utiliza-se
a diferença entre o log-retorno acumulado da ação e
o do Ibovespa:
\begin{equation}
s_{i,t}^{(N)}
=
\ln\left(\frac{P_{i,t}}{P_{i,t-N}}\right)
-
\ln\left(\frac{P_{m,t}}{P_{m,t-N}}\right),
\label{eq:spread_retorno_cap4}
\end{equation}
em que $P_{m,t}$ representa o nível do índice e $N$
é a janela retrospectiva.

As famílias de atributos contemplam diferentes janelas,
incluindo 3, 5, 14, 21, 42 e 63 pregões, além de medidas
de momentum em períodos maiores. Os grupos de atributos
são escolhidos por horizonte com dados da validação
de 2024.

Os atributos de entrada do XGBoost conjunto foram
organizados em três grupos cumulativos. O grupo básico
reúne retornos defasados, momentum, volatilidade, indicadores
de volume e de análise técnica, beta, retornos do Ibovespa
e do câmbio, além de posições relativas no ranking de
retornos das ações. O segundo grupo acrescenta medidas
intradiárias, indicadores de risco e retornos relativos
ao Ibovespa. O terceiro inclui também medidas do
comportamento conjunto das ações, como a mediana e a
dispersão dos retornos e a proporção de ações em alta.

O grupo básico foi utilizado no horizonte de 5 pregões,
o segundo no de 10 pregões e o terceiro no de 21 pregões.
Nos três horizontes, a identificação da ação também foi
incluída por meio de variáveis binárias.

A variável-alvo do XGBoost é o log-retorno acumulado
no horizonte $h$:
\begin{equation}
    y_{i,t}^{(h)}
    =
    \ln\left(\frac{P_{i,t+h}}{P_{i,t}}\right),
    \qquad h\in\{5,10,21\}.
    \label{eq:alvo_cap4}
\end{equation}

Esse retorno somente se torna conhecido no encerramento
do horizonte, em $t+h$. As etapas de treinamento e de
atualização da confiança respeitam essa data de
disponibilidade.

\subsection{Higienização e agregação das notícias}

As notícias passam por deduplicação e filtragem de
pertinência ao ativo ou à empresa. A agregação utilizada
pelo Black--Litterman considera a empresa emissora,
permitindo compartilhar informações entre diferentes
classes de ações.

Cada manchete recebe um escore
$s_k\in\{-1,0,+1\}$, correspondente às categorias
negativa, neutra e positiva. Para a empresa associada
à ação $i$, a média dos escores na data $d$ é
\begin{equation}
\bar{s}_{i,d}
=
\frac{1}{m_{i,d}}
\sum_{k=1}^{m_{i,d}}s_k,
\qquad m_{i,d}>0,
\label{eq:sent_medio_cap4}
\end{equation}
em que $m_{i,d}$ é a quantidade de manchetes elegíveis.

A ~\ref{tab:escores_noticias_trimestre} apresenta a
distribuição dos escores das manchetes por trimestre de
publicação. O escore médio trimestral é calculado pela
diferença entre o número de manchetes positivas e negativas,
dividida pelo total de manchetes no trimestre, incluindo
as neutras. As manchetes positivas foram a categoria mais
frequente em todos os trimestres do período.

\begin{table}[htbp]
    \centering
    \caption{Distribuição dos escores das manchetes por trimestre.}
    \label{tab:escores_noticias_trimestre}
    {\footnotesize
    \begin{tabular}{lrrrr}
        \hline
        Trimestre & Negativas & Neutras & Positivas & Escore médio \\
        \hline
        2023-T1 &  79 &  43 & 112 & 0,141 \\
        2023-T2 &  47 &  41 &  70 & 0,146 \\
        2023-T3 &  56 &  38 & 112 & 0,272 \\
        2023-T4 &  60 &  52 & 116 & 0,246 \\
        2024-T1 &  94 &  76 & 142 & 0,154 \\
        2024-T2 &  81 &  93 & 166 & 0,250 \\
        2024-T3 &  68 & 106 & 234 & 0,407 \\
        2024-T4 &  67 &  75 & 172 & 0,334 \\
        2025-T1 & 120 &  79 & 204 & 0,208 \\
        2025-T2 & 141 &  84 & 230 & 0,196 \\
        2025-T3 & 231 & 132 & 316 & 0,125 \\
        2025-T4 & 209 & 191 & 543 & 0,354 \\
        \hline
        Total & 1.253 & 1.010 & 2.417 & 0,249 \\
        \hline
    \end{tabular}}
    \par\smallskip
    Fonte: Elaborado pelo autor
\end{table}

O escore trimestral descreve a base textual. Para a
construção das opiniões do Black--Litterman, os escores
são agregados por empresa e por dia de publicação.

O sinal agregado é definido por
\begin{equation}
\psi_{i,d}
=
\tanh\left[
\bar{s}_{i,d}\ln(1+m_{i,d})
\right].
\label{eq:sent_suavizado_cap4}
\end{equation}

O logaritmo incorpora a quantidade de notícias com
ganhos decrescentes, enquanto a tangente hiperbólica
limita o sinal ao intervalo entre $-1$ e $1$.
Na ausência de notícias, o sinal é definido como zero.

Para cada decisão mensal, são consideradas somente
as notícias do dia civil anterior ao sinal. Essa data
pode corresponder a um dia sem pregão.

Como o sentimento é adimensional, sua conversão para
a escala de retorno utiliza a média dos valores
absolutos dos últimos dezoito retornos-alvo já
conhecidos, no horizonte selecionado para a ação.
A rotina exige pelo menos nove erros de previsão
maturados para construir o sinal e sua confiança.

\section{Modelo Utilizado}
\label{sec:modelos_utilizados}

Na simulação das carteiras, utiliza-se o XGBoost com treinamento conjunto das 20 ações para estimar os retornos incorporados às opiniões do modelo Black--Litterman. Em um experimento complementar, uma configuração do XGBoost treinada separadamente para cada ativo e horizonte é comparada a outros modelos de previsão de retorno, conforme descrito na \hyperref[sec:comparacao_modelos_previsao]{Seção~\ref*{sec:comparacao_modelos_previsao}}.

\subsection{Estimador quantitativo XGBoost}
\label{sec:estimador_xgb}

A configuração consolidada utiliza regressão XGBoost
conjunta das vinte ações, denominada \emph{pool}.
Para cada horizonte, os dados são reunidos e ordenados
por data e ativo. A identificação da ação é incluída
entre as variáveis de entrada por codificação binária.

São utilizados modelos por horizonte, com atualização
em janelas temporais sucessivas. A validação compreende
os dois semestres de 2024, e a avaliação utiliza os
quatro trimestres de 2025. Em cada etapa, o treinamento
reúne o histórico anterior à janela prevista e apenas
os alvos integralmente conhecidos antes de seu início.

O estimador quantitativo utiliza uma regressão XGBoost
com treinamento conjunto das 20 ações, abordagem
denominada \emph{pool}. Em cada janela de treinamento,
ajusta-se um modelo para cada horizonte de previsão
(5, 10 e 21 pregões), reunindo os registros das 20 ações.
Cada registro contém os atributos de uma ação em um
pregão e variáveis binárias que identificam essa ação.
Assim, o modelo pode aprender padrões comuns entre
as ações e produzir uma previsão de retorno individual
para cada uma. A ~\ref{tab:parametros_xgb}
apresenta a configuração utilizada nos três horizontes.

\begin{table}[H]
\centering
\caption{Configuração do XGBoost conjunto.}
\label{tab:parametros_xgb}
{\footnotesize
\renewcommand{\arraystretch}{0.9}
\begin{tabular}{p{0.55\linewidth}p{0.30\linewidth}}
\hline
Parâmetro & Valor \\ \hline
Função objetivo & Erro quadrático \\
Taxa de aprendizado & 0,03 \\
Profundidade máxima & 4 \\
Subamostragem de observações & 0,8 \\
Subamostragem de atributos & 0,8 \\
Peso mínimo por folha & 10 \\
Regularização L1 & 0,1 \\
Regularização L2 & 1,0 \\
Número máximo de árvores & 400 \\
Interrupção antecipada & 20 iterações sem melhora \\
Semente de aleatoriedade & 42 \\
\hline
\end{tabular}}
\par
Fonte: Elaborado pelo autor
\end{table}

O XGBoost conjunto utilizou uma configuração fixa de
referência, com erro quadrático e pesos iguais entre
as observações. Os valores da taxa de aprendizado, das
frações de subamostragem, da profundidade máxima, do
peso mínimo por folha e das penalizações L1 e L2 foram
retomados da configuração padrão do código de treinamento.
Esses valores foram mantidos nos três horizontes para
comparar os conjuntos de atributos e os métodos de
treinamento sob os mesmos hiperparâmetros.

Em cada janela, o número de árvores foi determinado
por validação temporal interna. Os primeiros 80\% dos
registros de treinamento foram destinados ao ajuste,
com exclusão das observações cujo retorno-alvo ainda
não era conhecido no início dos 20\% finais. Nessa
parcela final, monitorou-se o RMSE por até 400 iterações,
com interrupção após 20 iterações sem melhora.
Selecionado o número de árvores correspondente ao menor
RMSE, o modelo foi ajustado novamente com todos os
registros cujos retornos-alvo já eram conhecidos antes
do início da janela de previsão.

O XGBoost é treinado com log-retornos padronizados pela
média e pelo desvio-padrão dos dados usados no ajuste.
Antes de compor as opiniões do Black--Litterman, as
previsões voltam à escala original de log-retorno.

Para cada ação, escolheu-se um horizonte de 5, 10 ou
21 pregões com base no acerto direcional em 2024. A
escolha considerou apenas previsões cujos retornos já
eram conhecidos antes de 2025. Na simulação de 2025,
usou-se a previsão do XGBoost conjunto para o horizonte
escolhido de cada ação.

\subsection{Modelo de linguagem e classificação de sentimento}

A classificação das manchetes utiliza o modelo de
linguagem \texttt{gpt-5.6-luna}, em regime
\emph{zero-shot}, sem exemplos de classificação no
prompt. Em cada requisição, os campos
\texttt{\{ticker\}} e \texttt{\{headline\}} são
substituídos pelo código de negociação da ação e
pelo título da notícia, respectivamente.

\par\smallskip
\noindent
\begingroup
\setlength{\fboxsep}{5pt}
\fbox{%
    \begin{minipage}{\dimexpr\linewidth-2\fboxsep-2\fboxrule\relax}
        \small\raggedright
        \setlength{\parindent}{0pt}
        \setlength{\parskip}{0pt}

        \textbf{Prompt de classificação}\par
        Classifique o sentimento da manchete abaixo em relação
        ao impacto no preço da ação \texttt{\{ticker\}}.\par
        Responda apenas com uma palavra: positivo, negativo ou neutro.\par
        Manchete: \texttt{\{headline\}}
    \end{minipage}%
}
\endgroup
\par\smallskip

Para esta tarefa de análise de sentimento, não é realizado treinamento supervisionado adicional com as manchetes da amostra. Os rótulos são armazenados em arquivo estruturado e reaproveitados para pares de ativo e título já processados.

\subsection{Configuração do Black--Litterman}
\label{sec:configuracao_bl_cap4}

A integração utiliza expectativas diárias de retorno
excedente à referência livre de risco. A taxa anual
constante de 15\% é convertida para a frequência diária:
\begin{equation}
r_{f,d}
=
(1+0,15)^{1/252}-1.
\label{eq:selic_cap4}
\end{equation}

Em cada data de sinal, a matriz de covariância
$\Sigma_t$ é estimada com os retornos simples dos
últimos 126 pregões, incluindo uma regularização
diagonal de $10^{-7}$. A expectativa utilizada por
Markowitz corresponde à média desses retornos,
descontada de $r_{f,d}$.

O \emph{prior} do Black--Litterman é calculado por
\begin{equation}
    \pi_t
    =
    \delta\Sigma_t w^{\mathrm{eq}},
    \qquad
    w_i^{\mathrm{eq}}=\frac{1}{20},
    \qquad
    \delta=3.
    \label{eq:prior_experimental}
\end{equation}

Adotou-se $\delta=3$ como coeficiente fixo de aversão
ao risco, conforme a Subseção~\ref{subsec:otimizacao_pesos}.
O mesmo valor foi utilizado no cálculo dos retornos de
equilíbrio e na otimização das carteiras de Markowitz
e Black--Litterman.

A construção do \emph{prior} utiliza pesos de equilíbrio
iguais para as vinte ações do universo. Os pesos finais
de cada carteira são determinados pela otimização.

Seja $m_{i,t}$ a quantidade de notícias do dia elegível
para o sinal em $t$. A cobertura normalizada é
\begin{equation}
c_{i,t}
=
\min\left(1,\frac{m_{i,t}}{3}\right).
\label{eq:cobertura_cap4}
\end{equation}

A opinião híbrida no horizonte da ação é
\begin{equation}
Q_{i,t}^{(h)}
=
\left[
\alpha+(1-\alpha)(1-c_{i,t})
\right]\widehat{y}_{i,t}^{(h)}
+
(1-\alpha)c_{i,t}
\psi_{i,t}\rho_{i,t}^{(h)},
\label{eq:q_hibrido_cap4}
\end{equation}
em que $\widehat{y}_{i,t}^{(h)}$ é a previsão quantitativa,
$\rho_{i,t}^{(h)}$ é a escala de retorno do sentimento
e $\alpha=0,386251$.

A conversão para a escala diária excedente é
\begin{equation}
Q_{i,t}^{d}
=
\frac{Q_{i,t}^{(h)}}{h}
-
r_{f,d}.
\label{eq:q_daily_cap4}
\end{equation}

A divisão por $h$ produz uma taxa logarítmica diária
equivalente. Sua compatibilidade com os retornos
aritméticos do prior é aproximada.

\subsection{Controle de confiança das opiniões}

A confiança quantitativa utiliza os últimos dezoito
erros absolutos cujos retornos-alvo já foram realizados.
Seja $\mathcal{H}_{i,t}$ esse conjunto histórico e
$e_{i,u}$ o erro absoluto correspondente. Define-se
\begin{equation}
C_{i,t}^{\mathrm{modelo}}
=
\frac{1}{|\mathcal{H}_{i,t}|}
\sum_{u\in\mathcal{H}_{i,t}}
\exp\left(
-\frac{e_{i,u}}{\widetilde{e}_{i,t}}
\right),
\label{eq:confianca_modelo}
\end{equation}
em que $\widetilde{e}_{i,t}$ é a mediana desses erros,
limitada inferiormente por $10^{-8}$.

A confiança de notícias cresce com a quantidade
de manchetes:
\begin{equation}
C_{i,t}^{\mathrm{noticias}}
=
1-\exp\left(-\frac{m_{i,t}}{3}\right).
\label{eq:confianca_noticias}
\end{equation}

A combinação recebe pesos de 0,659122 para o componente
quantitativo e 0,340878 para o textual. O resultado
é limitado ao intervalo de 0,0575035 a 0,937247.

A variância diária equivalente da opinião é estimada
pela variância amostral dos retornos-alvo maturados,
dividida por $h^2$, com limite inferior de $10^{-8}$.
Denotando essa estimativa por $v_{i,t}^{d}$, a incerteza
básica é
\begin{equation}
\omega_{i,t}
=
\tau v_{i,t}^{d}
\frac{1-C_{i,t}}{C_{i,t}},
\qquad
\tau=0,0147235.
\label{eq:omega_cap4}
\end{equation}

Quando há menos de três manchetes, o mecanismo de
\emph{hard prior} reduz a confiança para $10^{-8}$ e
multiplica a incerteza resultante por 3,79598. Esse
aumento reduz a influência da opinião sobre o retorno
esperado, aproximando-o do retorno de equilíbrio.
O multiplicador foi mantido da configuração obtida
na calibração conjunta dos parâmetros do
Black--Litterman com o Optuna, que buscou valores
entre 1 e 6 na validação de 2024. A matriz $\Omega_t$
reúne essas incertezas na diagonal, com valores finais
limitados inferiormente por $10^{-8}$.

Como as opiniões são associadas diretamente às ações,
a atualização das expectativas pode ser expressa por
\begin{equation}
\mu_t^{\mathrm{BL}}
=
\pi_t
+
\tau\Sigma_t
\left(
\tau\Sigma_t+\Omega_t
\right)^{-1}
\left(
Q_t^{d}-\pi_t
\right).
\label{eq:posterior_experimental}
\end{equation}
\noindent onde, $\mu_t^{\mathrm{BL}}$ é o vetor de retornos
esperados atualizados, $\pi_t$ é o vetor de retornos de
equilíbrio (\emph{prior}) e $Q_t^{d}$ é o vetor de opiniões
em escala diária equivalente. A matriz $\Sigma_t$ representa
a covariância dos retornos, $\tau$ é o fator de escala da
incerteza do \emph{prior} e $\Omega_t$ reúne as incertezas
das opiniões. O termo $Q_t^{d}-\pi_t$ expressa o desvio das
opiniões em relação ao equilíbrio. Quanto maior a incerteza
de uma opinião, menor sua influência na atualização.

Segue o trecho completo atualizado. As tabelas têm uma única linha para o **XGBoost por ativo**, e o resultado da combinação XGBoost–LightGBM aparece em um parágrafo separado. Mantive as duas subseções e os rótulos existentes.

\subsection{Comparação complementar entre modelos de previsão}
\label{sec:comparacao_modelos_previsao}

Foi realizado um experimento complementar para comparar o desempenho de diferentes modelos na previsão de retorno. Foram utilizadas as mesmas 20 ações da base quantitativa, com o log-retorno futuro de fechamento como variável-alvo. A comparação considera os horizontes de 5, 10 e 21 pregões. Para cada horizonte, os modelos são avaliados nos mesmos pares de ativo e data. São incluídas apenas as previsões cujo horizonte termina até o último pregão de 2025, de modo que o retorno realizado possa ser calculado dentro do período de avaliação.

A comparação utiliza como referência as previsões do XGBoost conjunto descrito na \hyperref[sec:estimador_xgb]{Seção~\ref*{sec:estimador_xgb}}, que foram incorporadas às opiniões do modelo Black--Litterman na simulação das carteiras. Para cada horizonte e trimestre, o LightGBM, a regressão Ridge e a rede recorrente do tipo GRU são ajustados com os dados reunidos das 20 ações. Assim, cada modelo é treinado conjuntamente e produz previsões individuais para cada ativo.

O LightGBM e a regressão Ridge utilizam os mesmos atributos tabulares do XGBoost. A regressão Ridge constitui uma referência linear com regularização L2. A GRU utiliza sequências de atributos dos últimos 35 pregões para estimar o retorno futuro. Na configuração denominada Lag-Llama com Ridge, os históricos de log-retornos de cada ativo são utilizados para extrair representações do Lag-Llama pré-treinado. Essas representações são utilizadas como entrada de uma regressão Ridge ajustada conjuntamente com os dados das 20 ações.

A avaliação é realizada nos quatro trimestres de 2025. Os modelos são ajustados no início de cada trimestre com os dados das janelas anteriores, utilizando apenas observações cujos retornos-alvo já eram conhecidos. Além da avaliação por horizonte, considera-se um recorte que reúne, para cada ação, as previsões no horizonte adotado na simulação das carteiras. Esse recorte é denominado Horizontes do TCC nas tabelas de resultados.

Para interpretar a acurácia direcional, considera-se a proporção de retornos positivos, que corresponde à taxa de acerto obtida ao prever alta em todas as observações. As métricas utilizadas são apresentadas na \hyperref[sec:desempenho_das_previsoes_quantitativas]{Seção~\ref*{sec:desempenho_das_previsoes_quantitativas}}.

\subsection{Resultados da comparação entre modelos}
\label{sec:resultados_comparacao_modelos}

A avaliação compreende 4.900, 4.800 e 4.580 pares de ativo e data para os horizontes de 5, 10 e 21 pregões, respectivamente. O número de observações diminui à medida que o horizonte aumenta, pois uma parcela maior das previsões realizadas no final do ano tem seu horizonte encerrado após o período de avaliação. Também são avaliadas 4.683 observações no recorte denominado Horizontes do TCC, que utiliza, para cada ação, o horizonte adotado na simulação das carteiras.

A Tabela~\ref{tab:comparacao_modelos_rmse} apresenta o MAE e o RMSE calculados sobre todas as observações de cada recorte. A linha Lag-Llama com Ridge corresponde à configuração em que as representações do modelo pré-treinado são utilizadas em uma regressão ajustada conjuntamente com os dados das 20 ações.

\begin{table}[htbp]
    \centering
    \caption{Erros das previsões de retorno por modelo e horizonte em 2025 (\%).}
    \label{tab:comparacao_modelos_rmse}
    {\small
    \setlength{\tabcolsep}{4pt}
    \begin{tabular}{@{}lrrrrrrrr@{}}
        \hline
        & \multicolumn{2}{c}{5 pregões}
        & \multicolumn{2}{c}{10 pregões}
        & \multicolumn{2}{c}{21 pregões}
        & \multicolumn{2}{c}{Horizontes do TCC} \\
        \cline{2-9}
        Modelo & MAE & RMSE & MAE & RMSE & MAE & RMSE & MAE & RMSE \\
        \hline
        XGBoost             & 3,099 & \textbf{4,309} & 4,539 & \textbf{6,177} & 6,719 & 9,085 & 5,654 & 8,082 \\
        LightGBM            & \textbf{3,098} & 4,318 & \textbf{4,533} & 6,177 & \textbf{6,714} & \textbf{9,077} & \textbf{5,648} & \textbf{8,070} \\
        GRU                 & 3,139 & 4,355 & 4,591 & 6,245 & 7,004 & 9,388 & 5,841 & 8,279 \\
        Lag-Llama com Ridge & 3,176 & 4,384 & 4,674 & 6,378 & 7,103 & 9,615 & 5,986 & 8,527 \\
        Ridge               & 3,147 & 4,366 & 4,706 & 6,408 & 7,449 & 9,946 & 6,176 & 8,742 \\
        \hline
    \end{tabular}
    }
    \par\smallskip
    Fonte: Elaborado pelo autor
\end{table}

O XGBoost conjunto apresentou MAE e RMSE menores que os da GRU, do Lag-Llama com Ridge e da Ridge conjunta nos três horizontes. Os erros do XGBoost e do LightGBM foram próximos. No recorte dos horizontes do TCC, o LightGBM apresentou RMSE de 8,070\%, frente a 8,082\% do XGBoost conjunto. A diferença de RMSE entre LightGBM e XGBoost foi de $-0{,}012$ ponto percentual, com intervalo de confiança de 95\% de $-0{,}022$ a $-0{,}002$ ponto percentual.

A ~\ref{tab:comparacao_modelos_direcao} apresenta a acurácia direcional por horizonte e a acurácia balanceada no recorte dos horizontes do TCC. A acurácia balanceada atribui o mesmo peso às classes de retorno positivo e não positivo.

\begin{table}[htbp]
    \centering
    \caption{Acurácia direcional e balanceada das previsões de retorno em 2025 (\%).}
    \label{tab:comparacao_modelos_direcao}
    {\small
    \setlength{\tabcolsep}{4pt}
    \begin{tabular}{@{}lrrrrr@{}}
        \hline
        & \multicolumn{3}{c}{Acurácia direcional por horizonte}
        & \multicolumn{2}{c}{Horizontes do TCC} \\
        \cline{2-4}\cline{5-6}
        Modelo & 5 pregões & 10 pregões & 21 pregões & Direcional & Balanceada \\
        \hline
        XGBoost             & \textbf{53,92} & 56,21 & 58,41 & \textbf{56,46} & \textbf{50,25} \\
        LightGBM            & 53,14 & \textbf{56,67} & \textbf{58,43} & 56,18 & 50,03 \\
        GRU                 & 49,55 & 51,12 & 45,59 & 47,23 & 46,03 \\
        Lag-Llama com Ridge & 49,78 & 51,06 & 50,04 & 48,49 & 48,01 \\
        Ridge               & 50,69 & 49,96 & 47,03 & 46,57 & 45,94 \\
        \hline
    \end{tabular}
    }
    \par\smallskip
    Fonte: Elaborado pelo autor
\end{table}

Nos horizontes de 10 e 21 pregões, o LightGBM apresentou a maior acurácia direcional, com valores próximos aos do XGBoost conjunto. No recorte dos horizontes do TCC, o XGBoost apresentou acurácia direcional de 56,46\%, frente a 56,18\% do LightGBM. A diferença entre LightGBM e XGBoost foi de $-0{,}28$ ponto percentual, com intervalo de confiança de 95\% de $-0{,}80$ a $0{,}15$ ponto percentual. O intervalo inclui zero, indicando que a diferença observada é inconclusiva.

Nesse mesmo recorte, 56,44\% dos retornos realizados foram positivos, valor que corresponde à acurácia direcional obtida ao prever alta em todas as observações. O XGBoost conjunto apresentou acurácia balanceada de 50,25\%, frente a 50,00\% da regra de sempre prever alta. Assim, a acurácia direcional observada deve ser interpretada à luz da proporção de altas no período.

Na comparação de 2025, o XGBoost conjunto apresentou erros menores e maior acurácia direcional que a GRU, o Lag-Llama com Ridge e a Ridge conjunta. O LightGBM apresentou desempenho próximo ao do XGBoost nas duas dimensões, com pequena vantagem no erro de magnitude no recorte dos horizontes do TCC. Esses resultados se referem à previsão dos retornos. O desempenho financeiro das carteiras é avaliado na simulação Black--Litterman com o XGBoost conjunto.

\section{Protocolo de Alocação e Simulação}
\label{sec:protocolo_experimental}

\subsection{Restrições e otimização}

O capital inicial de R\$~100.000 é utilizado como
referência para apresentar a evolução do patrimônio e
os custos de negociação em reais. Como os pesos e os
custos de negociação são proporcionais ao capital, seu
valor inicial não altera os pesos ótimos nem os
retornos percentuais da carteira.

A simulação exige pesos positivos para todas as ações
selecionadas, o que impede posições vendidas. A soma
dos pesos é igual a~1, de modo que todo o capital
disponível é alocado sem recorrer a empréstimos.
Para uma carteira com $n$ ações, os pesos são obtidos
pela maximização de
\begin{equation}
\max_w
\left[
\mu_t^\top w
-
\frac{\delta}{2}
w^\top\Sigma_t w
\right],
\label{eq:otimizacao_experimental}
\end{equation}
sujeita a
\begin{equation}
\sum_{i=1}^{n}w_i=1,
\qquad
0,001\leq w_i\leq
\max\left(0,341677,\frac{1}{n}\right).
\label{eq:restricoes_experimentais}
\end{equation}

O piso mantém uma participação mínima de 0,1\% por
ação selecionada. O teto se ajusta à quantidade de
ativos para permitir carteiras com uma ou duas ações.
A aversão ao risco é fixada em $\delta=3$.

\subsection{Ativos comuns e seleção própria}

A partir de um universo fixo de 20 ações, a simulação
avalia carteiras com $n$ ações, de 1 a 20, em dois
cenários: ativos comuns e seleção própria.

No cenário de ativos comuns, cada grupo é um subconjunto
de $n$ ações do universo. O mesmo grupo é utilizado por
todas as configurações comparadas e permanece fixo durante
o ano. A cada rebalanceamento, são recalculados apenas os
pesos, o que permite comparar as configurações sob a mesma
composição de ativos.

Para cada tamanho de carteira entre 2 e 18 ações,
sorteiam-se 30 combinações distintas de $n$ ações entre
as possíveis no universo de 20. Por exemplo, para $n=2$,
sorteiam-se 30 dos 190 pares possíveis. Dessa forma,
um mesmo tamanho de carteira é avaliado com diferentes
composições.

Para $n=1$ e $n=19$, são avaliadas as 20 combinações
possíveis de cada tamanho. Para $n=20$, avalia-se apenas
a carteira que reúne todas as ações. Ao todo, o cenário
de ativos comuns inclui 551 combinações por configuração.
O sorteio utiliza a mesma semente aleatória para permitir a
reprodução das combinações selecionadas.

No cenário de seleção própria, cada configuração
seleciona mensalmente as $n$ ações com maior pontuação,
calculada por
\begin{equation}
    a_{i,t}
    =
    \frac{\mu_{i,t}}{\sqrt{\Sigma_{ii,t}}},
    \label{eq:ranking_selecao}
\end{equation}
em que $\mu_{i,t}$ é o retorno excedente esperado da
ação $i$ na data $t$ e $\sqrt{\Sigma_{ii,t}}$ é a
volatilidade estimada dessa ação. Após a seleção,
os pesos são calculados pelo mesmo otimizador. Tanto
a composição quanto os pesos da carteira podem mudar
a cada rebalanceamento.

A comparação entre Black--Litterman e Markowitz totaliza
1.142 simulações. Para cada estratégia, são realizadas
551 simulações com ativos comuns e 20 com seleção própria,
uma para cada valor de $n$. Os controles com \emph{prior}
sem opiniões e XGBoost direto seguem o mesmo desenho e
somam outras 1.142 simulações. A carteira de pesos iguais
é avaliada no cenário de ativos comuns, acrescentando
uma simulação para cada um dos 551 grupos de ações.

O calendário de rebalanceamento é o mesmo para todas
as configurações em cada cenário. A frequência mensal permite
atualizar as carteiras com as informações disponíveis a
cada mês, com menos ocasiões de negociação do que em um
rebalanceamento diário. O sinal de alocação é calculado
no fechamento do primeiro pregão de cada mês, e a nova
alocação é executada no fechamento do pregão seguinte.

Até a primeira compra de ações, o dinheiro fica em caixa
e rende à taxa livre de risco adotada na simulação. Nos
rebalanceamentos seguintes, calcula-se primeiro o valor
da carteira existente com os preços de fechamento do
dia. Em seguida, as compras e vendas são feitas a esses
mesmos preços. A nova carteira começa a render a partir
do pregão seguinte.

Entre rebalanceamentos, a quantidade de ações de cada
empresa permanece a mesma. A participação de cada ação
no valor da carteira, porém, muda com os preços. Por
exemplo, se uma ação sobe enquanto as demais mantêm
seus preços, ela passa a representar uma parcela maior
da carteira.

\subsection{Custos e evolução do patrimônio}

O custo de negociação é calculado multiplicando o
patrimônio antes da operação pelo turnover e pela
taxa de 0,1\%. O turnover considera as mudanças
nos valores investidos em cada ação e no caixa.

Sejam $C_t^{-}$ o patrimônio imediatamente antes
da negociação, $H_{i,t}^{-}$ o valor de cada posição
existente e $B_t^{-}$ o caixa disponível. Para os
pesos-alvo $w_{i,t}$, o turnover é
\begin{equation}
\mathrm{TO}_t
=
\frac{
\displaystyle
\sum_i
\left|
w_{i,t}C_t^{+}-H_{i,t}^{-}
\right|
+
|B_t^{-}|
}{
2C_t^{-}
},
\label{eq:turnover_cap4}
\end{equation}
em que $C_t^{+}$ é o patrimônio após o pagamento
dos custos:
\begin{equation}
C_t^{+}
=
C_t^{-}
\left(
1-0,001\,\mathrm{TO}_t
\right).
\label{eq:custo_autofinanciado}
\end{equation}

As duas expressões são resolvidas conjuntamente,
pois o custo reduz o valor disponível para formar
as novas posições. Esse procedimento incorpora
a movimentação efetiva e o financiamento dos custos
pelo próprio patrimônio.

O retorno líquido diário e o retorno acumulado são
calculados por
\begin{equation}
R_t^{p}
=
\frac{C_t}{C_{t-1}}-1,
\qquad
R_{\mathrm{acum}}
=
\frac{C_T}{C_0}-1.
\label{eq:ret_acumulado_cap4}
\end{equation}

O cenário sem custos mantém os pesos-alvo e o
calendário de operações, retirando apenas a cobrança
financeira. Não há nova otimização nesse cenário.
A simulação não inclui liquidação terminal das
posições.

\subsection{Configurações de comparação}

A ~\ref{tab:configuracoes_alocacao} apresenta as
configurações avaliadas na simulação financeira.

\begin{table}[H]
    \centering
    \caption{Configurações de alocação avaliadas.}
    \label{tab:configuracoes_alocacao}
    {\footnotesize
    \renewcommand{\arraystretch}{1.25}
    \setlength{\tabcolsep}{5pt}
    \begin{tabular}{@{}p{0.22\linewidth}p{0.74\linewidth}@{}}
        \hline
        \textbf{Configuração} &
        \textbf{Retorno esperado ou regra de alocação}
        \tabularnewline
        \hline

        BL híbrido &
        \raggedright Retorno atualizado a partir do \emph{prior},
        da previsão do XGBoost e do sentimento das manchetes.
        \tabularnewline

        Markowitz &
        \raggedright Média dos retornos simples dos últimos
        126 pregões, menos a taxa livre de risco diária.
        \tabularnewline

        \emph{Prior} sem opiniões &
        \raggedright Retorno de equilíbrio $\pi_t$ usado
        diretamente na otimização.
        \tabularnewline

        XGBoost direto &
        \raggedright Previsão do XGBoost dividida pelo horizonte,
        menos a taxa livre de risco diária.
        \tabularnewline

        Pesos iguais ($1/n$) &
        \raggedright Em cada rebalanceamento, cada ação recebe
        $1/n$ do capital.
        \tabularnewline

        \hline
    \end{tabular}}
    \par\smallskip
    Fonte: Elaborado pelo autor
\end{table}

As quatro configurações que calculam os pesos por
otimização utilizam a mesma matriz de covariância,
o mesmo coeficiente de aversão ao risco e as mesmas
restrições. A carteira de pesos iguais usa os mesmos
grupos de ações, o mesmo calendário mensal e a mesma
regra de custos. Como seus pesos são definidos diretamente
por $w_i=1/n$, ela serve de referência para avaliar o
efeito da otimização dos pesos em carteiras com as mesmas
ações. A referência de renda fixa corresponde à
capitalização da taxa constante de 15\% ao ano.
\newpage
A ~\ref{tab:parametros_vigentes_cap4}
reúne as principais condições da simulação.

\begin{table}[H]
\centering
\caption{Parâmetros da simulação financeira consolidada.}
\label{tab:parametros_vigentes_cap4}
{\footnotesize
\renewcommand{\arraystretch}{0.9}
\begin{tabular}{p{0.55\linewidth}p{0.32\linewidth}}
\hline
Parâmetro & Valor \\ \hline
Capital inicial & R\$ 100.000 \\
Quantidade de ações & De 1 a 20 \\
Janela de estimação & 126 pregões \\
Aversão ao risco & 3 \\
Regularização da covariância & $10^{-7}$ \\
Pesos de equilíbrio do prior & Iguais entre as vinte ações \\
Peso quantitativo, $\alpha$ & 0,386251 \\
Escala do prior, $\tau$ & 0,0147235 \\
Gate de notícias & 3 manchetes \\
Histórico de erros maturados & Últimas 18 observações \\
Mínimo de erros maturados & 9 observações \\
Piso por ação & 0,1\% \\
Teto por ação & $\max(34,1677\%;1/n)$ \\
Rebalanceamento & Mensal \\
Execução & Fechamento seguinte ao sinal \\
Custo de negociação & 0,1\% sobre o turnover \\
Referência livre de risco & 15\% ao ano, constante \\
\hline
\end{tabular}}
\par
Fonte: Elaborado pelo autor
\end{table}

\section{Métricas de Avaliação}
\label{sec:metricas_experimentais}

\subsection{Desempenho das previsões quantitativas}
\label{sec:desempenho_das_previsoes_quantitativas}

A acurácia direcional, ou \emph{hit ratio} (HR),
mede a proporção de previsões cujo sinal coincide
com o retorno realizado:
\begin{equation}
\mathrm{HR}
=
\frac{1}{J}
\sum_{j=1}^{J}
\mathbf{1}
\left\{
\operatorname{sgn}(\widehat{y}_j)
=
\operatorname{sgn}(y_j)
\right\}.
\label{eq:hr_cap4}
\end{equation}

A comparação inclui a acurácia majoritária,
correspondente à frequência da classe de retorno
mais comum entre positiva, nula e negativa.
Essa referência indica o acerto obtido ao prever
sempre a mesma classe.

O erro absoluto médio e a raiz do erro quadrático
médio são definidos por
\begin{equation}
\mathrm{MAE}
=
\frac{1}{J}
\sum_{j=1}^{J}
|\widehat{y}_j-y_j|,
\label{eq:mae_cap4}
\end{equation}
\begin{equation}
\mathrm{RMSE}
=
\sqrt{
\frac{1}{J}
\sum_{j=1}^{J}
(\widehat{y}_j-y_j)^2
}.
\label{eq:rmse_cap4}
\end{equation}

O coeficiente de determinação é
\begin{equation}
R^2
=
1-
\frac{
\sum_{j=1}^{J}(y_j-\widehat{y}_j)^2
}{
\sum_{j=1}^{J}(y_j-\bar{y})^2
},
\label{eq:r2_cap4}
\end{equation}
em que $\bar{y}$ é a média dos retornos realizados
na amostra avaliada.

As métricas são calculadas por ativo e horizonte.
Os resumos entre ações utilizam médias simples
dos indicadores individuais. Quando todas as
observações são reunidas antes do cálculo,
o resultado é identificado como agregado.

A avaliação de erros em 2025 considera somente
previsões com retorno-alvo disponível até o
encerramento do período. Assim, cada ação possui
245, 240 ou 229 pares de previsão e realização,
respectivamente, para os horizontes de 5, 10
e 21 pregões. A simulação financeira utiliza
também previsões disponíveis para decisão cujo
resultado completo ainda não era conhecido
ao final do ano.

\subsection{Avaliação da classificação de sentimento}

A classificação textual é comparada com uma
referência estabelecida por um único avaliador
humano em uma amostra aleatória simples de
300 manchetes selecionadas aleatoriamente.
A definição dos rótulos conta com apoio de IA
na discussão de casos ambíguos.

São calculadas a concordância total, a precisão,
a revocação e o F1 de cada classe. Para uma
classe $c$, essas medidas são
\begin{equation}
\mathrm{Precisao}_c
=
\frac{\mathrm{VP}_c}
{\mathrm{VP}_c+\mathrm{FP}_c},
\qquad
\mathrm{Revocacao}_c
=
\frac{\mathrm{VP}_c}
{\mathrm{VP}_c+\mathrm{FN}_c},
\label{eq:precisao_revocacao}
\end{equation}
\begin{equation}
\mathrm{F1}_c
=
\frac{
2\,\mathrm{Precisao}_c\,\mathrm{Revocacao}_c
}{
\mathrm{Precisao}_c+\mathrm{Revocacao}_c
},
\label{eq:f1_sentimento}
\end{equation}
em que VP, FP e FN representam verdadeiros
positivos, falsos positivos e falsos negativos
da classe considerada.

O macro-F1 corresponde à média dos valores de F1
das três classes. O kappa de Cohen desconta
a concordância esperada pelas frequências
das classificações:
\begin{equation}
\kappa
=
\frac{p_o-p_e}{1-p_e},
\label{eq:kappa_sentimento}
\end{equation}
em que $p_o$ é a concordância observada e $p_e$
é a concordância esperada pelas frequências
marginais.

Os rótulos humanos são utilizados na avaliação
da classificação e não substituem os rótulos
das notícias utilizados na simulação financeira.

\subsection{Avaliação preditiva das camadas intermediárias}

A análise complementar das camadas utiliza
as previsões e os sinais da execução de referência
registrada para essa finalidade, apresentada
separadamente da rodada financeira consolidada
com o modelo conjunto.

São comparados XGBoost, Q completo, Q sem polaridade
textual, BL com XGBoost, BL sem polaridade,
BL completo e prior sem views. Nos controles
sem polaridade, o sentimento é fixado em zero,
preservando a cobertura, os pesos quantitativos,
a confiança, o gate e a covariância. No BL com
XGBoost, a opinião recebe apenas a previsão
quantitativa, mantendo a matriz de incerteza
da referência.

O alvo principal é o retorno excedente diário
equivalente:
\begin{equation}
z_{i,t}^{(h)}
=
\frac{1}{h}
\ln\left(
\frac{P_{i,t+h}}{P_{i,t}}
\right)
-
r_{f,d}.
\label{eq:alvo_diario_camadas}
\end{equation}

A avaliação utiliza 228 pares ativo-data com
resultado disponível no horizonte selecionado.
Também é realizada uma avaliação contra o
retorno simples excedente do pregão seguinte:
\begin{equation}
z_{i,t}^{(1)}
=
\frac{P_{i,t+1}}{P_{i,t}}
-1-r_{f,d},
\label{eq:alvo_pregao_seguinte}
\end{equation}
com os 240 pares das decisões mensais.

Na avaliação do pregão seguinte, as previsões
logarítmicas de XGBoost e Q são convertidas
para retornos simples por exponenciação,
com o ajuste da referência livre de risco.
Prior e posterior mantêm sua escala aritmética
diária excedente.

Além de HR, MAE, RMSE e $R^2$, utiliza-se
a acurácia balanceada:
\begin{equation}
\mathrm{BA}
=
\frac{1}{K}
\sum_{c=1}^{K}
\frac{
\text{acertos na classe }c
}{
\text{observacoes da classe }c
},
\label{eq:acuracia_balanceada}
\end{equation}
em que $K$ é a quantidade de classes de sinal
presentes na amostra.

As diferenças entre configurações são calculadas
sobre as mesmas observações. A redução relativa
do erro quadrático ao passar da configuração
$A$ para $B$ é
\begin{equation}
\mathrm{ReducaoMSE}_{A\rightarrow B}
=
100\left(
1-\frac{\mathrm{MSE}_B}{\mathrm{MSE}_A}
\right).
\label{eq:reducao_mse_camadas}
\end{equation}

Valores positivos indicam redução do erro.
As diferenças de HR são apresentadas em pontos
percentuais. Os erros diários equivalentes
são expressos em pontos-base, sendo
100 pontos-base iguais a um ponto percentual.

\subsection{Desempenho financeiro e risco}

O retorno acumulado é acompanhado pela
volatilidade anualizada dos retornos líquidos:
\begin{equation}
\sigma_{\mathrm{anual}}
=
s(R^p)\sqrt{252},
\label{eq:volatilidade_anual}
\end{equation}
em que $s(R^p)$ é o desvio-padrão amostral
dos retornos diários da carteira.

O índice de Sharpe utiliza o retorno excedente
à referência livre de risco:
\begin{equation}
\mathrm{Sharpe}
=
\frac{
\overline{R^p-r_{f,d}}
}{
s(R^p-r_{f,d})
}
\sqrt{252}.
\label{eq:sharpe_cap4}
\end{equation}

O drawdown representa a queda do patrimônio
em relação ao maior valor anteriormente atingido:
\begin{equation}
\mathrm{DD}_t
=
\frac{
C_t
}{
\max_{0\leq u\leq t}C_u
}
-1,
\qquad
\mathrm{MDD}
=
\min_t\mathrm{DD}_t.
\label{eq:drawdown_cap4}
\end{equation}

A avaliação também registra o turnover acumulado,
os custos pagos e a diferença entre os retornos
com e sem custos.

\subsection{Comparação entre grupos de ações}

No cenário de ativos comuns, os resultados
são resumidos por mediana e percentis para
cada quantidade de ações. Cada indicador
é calculado separadamente entre os grupos
avaliados.

A diferença de retorno entre os modelos
é calculada primeiro em cada grupo:
\begin{equation}
\Delta_j
=
R_j^{\mathrm{BL}}
-
R_j^{\mathrm{MV}}.
\label{eq:diferenca_grupo}
\end{equation}

Em seguida, apresenta-se a mediana dessas
diferenças, denominada $\Delta$ mediano.
Também é calculado o percentual de grupos
em que o retorno do BL supera estritamente
o de Markowitz. Os empates não são contabilizados
como vitórias.

A carteira de pesos iguais é comparada ao BL híbrido
e a Markowitz grupo a grupo, com base no retorno
líquido. Para cada comparação, calcula-se a mediana
das diferenças de retorno e a proporção de grupos
em que cada estratégia obteve o maior retorno.
A análise também considera separadamente os grupos
com três a vinte ações, pois, com uma ou duas ações,
as restrições impõem os mesmos pesos no cenário
de ativos comuns.

Na seleção própria, são apresentados diretamente
os indicadores da simulação de cada configuração
e quantidade de ações, pois existe uma carteira
dinâmica por combinação avaliada.